\begin{scanpage}{0086}{69}
\[
\left.\begin{aligned}
E_\vartheta=H_\varphi&=\frac{iaH_{0x}}2\{[g^1(1)-g^1(2)]J_1(\xi)\\
&\qquad+i[g^1(1)+g^1(2)]J_2(\xi)\}\frac{e^{ikR}}R.
\end{aligned}\right\}\tag{12.06}
\]
where
\[\left.\begin{aligned}\zeta&=ka(\sin\vartheta-\sin\gamma),\\\xi&=ka(\sin\vartheta+\sin\gamma).\end{aligned}\right\}\tag{12.07}\]
These expressions are valid in the region $0<\vartheta<\frac\pi2$; when $|\zeta|\gg1$ and $|\xi|\gg1$ they change to Equations (12.01) and (12.02), and in the direction $\vartheta=\gamma$, $\varphi=\frac\pi2$ they give a field equal to zero.

Using specific expressions for the functions $f^1$ and $g^1$, it is not difficult to establish that the total field scattered by the disk in view of Equations (10.03) and (10.04) may be represented in the following form:

with $\varphi=\frac\pi2$
\[
\left.\begin{aligned}
E_\varphi=-H_\vartheta&=\frac{iaE_{0x}}2\{[f(1)-f(2)]J_1(\zeta)\\
&\qquad-i[f(1)+f(2)]J_2(\zeta)\}\frac{e^{ikR}}R,\\
E_\vartheta=H_\varphi&=\frac{iaH_{0x}}2\{[g(1)-g(2)]J_1(\zeta)\\
&\qquad-i[g(1)+g(2)]J_2(\zeta)\}\frac{e^{ikR}}R,
\end{aligned}\right\}\tag{12.08}
\]
and with $\varphi=-\frac\pi2$
\[
\left.\begin{aligned}
E_\varphi=-H_\vartheta&=\frac{iaE_{0x}}2\{[f(1)-f(2)]J_1(\xi)\\
&\qquad+i[f(1)+f(2)]J_2(\xi)\}\frac{e^{ikR}}R,\\
E_\vartheta=H_\varphi&=\frac{iaH_{0x}}2\{[g(1)-g(2)]J_1(\xi)\\
&\qquad+i[g(1)+g(2)]J_2(\xi)\}\frac{e^{ikR}}R.
\end{aligned}\right\}\tag{12.09}
\]
It is convenient to write these expressions as follows
\[
\left.\begin{aligned}
E_\varphi=-H_\vartheta&=\frac{iaE_{0x}}2\overline\Sigma(\vartheta,\gamma)\frac{e^{ikR}}R,\\
E_\vartheta=H_\varphi&=\frac{iaH_{0x}}2\Sigma(\vartheta,\gamma)\frac{e^{ikR}}R,
\end{aligned}\right\}\tag{12.10}
\]
\end{scanpage}
